\documentclass[12pt,a4paper]{article}

% -----------------------------------------------------------------------
% Packages
% -----------------------------------------------------------------------
\usepackage[utf8]{inputenc}
\usepackage[T1]{fontenc}
\usepackage{geometry}
\usepackage{hyperref}
\usepackage{listings}
\usepackage{xcolor}
\usepackage{graphicx}
\usepackage{enumitem}
\usepackage{titlesec}
\usepackage{parskip}
\usepackage{booktabs}
\usepackage{array}
\usepackage{fancyhdr}
\usepackage{amsmath}

\geometry{margin=2.5cm}

% -----------------------------------------------------------------------
% Hyperref setup
% -----------------------------------------------------------------------
\hypersetup{
    colorlinks=true,
    linkcolor=black,
    urlcolor=blue,
    citecolor=black
}

% -----------------------------------------------------------------------
% Code listing style
% -----------------------------------------------------------------------
\definecolor{codebg}{rgb}{0.95,0.95,0.97}
\definecolor{codecomment}{rgb}{0.4,0.4,0.6}
\definecolor{codestring}{rgb}{0.13,0.55,0.13}
\definecolor{codekeyword}{rgb}{0.0,0.35,0.75}

\lstdefinestyle{pythonstyle}{
    backgroundcolor=\color{codebg},
    commentstyle=\color{codecomment}\itshape,
    keywordstyle=\color{codekeyword}\bfseries,
    stringstyle=\color{codestring},
    basicstyle=\ttfamily\small,
    breaklines=true,
    keepspaces=true,
    numbers=left,
    numberstyle=\tiny\color{gray},
    showspaces=false,
    showstringspaces=false,
    tabsize=4,
    frame=single,
    framesep=4pt,
    rulecolor=\color{gray!40},
}

\lstdefinestyle{bashstyle}{
    backgroundcolor=\color{codebg},
    basicstyle=\ttfamily\small,
    breaklines=true,
    frame=single,
    framesep=4pt,
    rulecolor=\color{gray!40},
    numbers=none,
}

% -----------------------------------------------------------------------
% Header / footer
% -----------------------------------------------------------------------
\pagestyle{fancy}
\fancyhf{}
\rhead{Mini Game Hub}
\lhead{CS108 Project Report}
\cfoot{\thepage}
\renewcommand{\headrulewidth}{0.4pt}

% -----------------------------------------------------------------------
% Title
% -----------------------------------------------------------------------
\title{
    \vspace{-1cm}
    {\Large CS108 Project Report} \\[0.4em]
    {\huge \textbf{Mini Game Hub}} \\[0.4em]
    {\large Bash $\cdot$ Python $\cdot$ Pygame}
}
\author{Lucky Gupta and Duduku Jaswanth Balaji}
\date{April 2025}

% -----------------------------------------------------------------------
% Document
% -----------------------------------------------------------------------
\begin{document}

\maketitle
\thispagestyle{fancy}
\tableofcontents
\newpage

% ===================================================================
\section{Introduction}
% ===================================================================

Mini Game Hub is a Two-player game hub built as part of the CS108 course
(Spring semester 2026). It integrates Bash scripting for user authentication and
Python (Pygame) for graphical gameplay, enabling two authenticated players to
select a game from an interactive menu, play via a Pygame GUI window, and have
their results recorded on a persistent leaderboard.

The project focuses on modularity, code quality, and a clean user experience.
It supports the following features:

\begin{itemize}
    \item Two-player authentication with SHA-256 hashed passwords stored in
          \texttt{users.tsv}.
    \item An interactive Pygame game-selection menu with a single-click launch.
    \item Three fully playable board games rendered with Pygame:
          Tic-Tac-Toe (10$\times$10), Othello (8$\times$8), and Connect Four
          (7$\times$7).
    \item Persistent game history recorded in \texttt{history.csv} after every
          game.
    \item A formatted terminal leaderboard rendered by a Bash script, sortable
          by wins, losses, or win/loss ratio.
    \item A post-game screen that prompts players to select a leaderboard sort
          metric before displaying statistics.
\end{itemize}

% ===================================================================
\section{Running Instructions}
% ===================================================================

\begin{itemize}
    \item Ensure Python 3.10+ and \texttt{pip} are installed.
    \item Install the required Python libraries:
\begin{lstlisting}[style=bashstyle]
pip install pygame-ce numpy matplotlib
\end{lstlisting}
    \item Navigate to the project root directory and run:
\begin{lstlisting}[style=bashstyle]
bash main.sh
\end{lstlisting}
    \item Follow the authentication prompts for both players, then select a
          game from the Pygame menu that appears.
    
\end{itemize}

% ===================================================================
\section{Project Layout}
% ===================================================================

\subsection{Authentication (\texttt{main.sh})}

\texttt{main.sh} is the entry point for the hub. It authenticates two players
before launching the Python game engine.

\begin{itemize}
    \item Prompts for two distinct usernames and passwords sequentially.
    \item Prevents both players from using the same username.
    \item Hashes passwords with SHA-256 using \texttt{sha256sum} before storing
          them in \texttt{users.tsv} --- plaintext passwords are never stored.
    \item On login, hashes the typed password and compares it against the
          stored hash; re-prompts on mismatch.
    \item If a username does not exist, offers registration and updates
          \texttt{users.tsv} accordingly.
    \item Once both players are authenticated, calls:
\begin{lstlisting}[style=bashstyle]
python3 game.py <username1> <username2>
\end{lstlisting}
\end{itemize}

\subsection{Game Engine (\texttt{game.py})}

\texttt{game.py} receives both authenticated usernames and manages the full
Python-side flow.

\begin{itemize}
    \item Initialises a 700$\times$800 Pygame window.
    \item Displays the game-selection \textbf{MenuScreen}: a plain vertical
          list of games, single click to launch --- no hover card effects, no
          confirmation button.
    \item Runs the selected game via \texttt{run\_game()}, which manages the
          Pygame event loop generically across all three games.
    \item After each game, displays the \textbf{PostGameScreen} in two phases:
          \begin{enumerate}
              \item Shows the result and three sort-metric buttons
                    (Wins / Losses / W/L Ratio). On selection, writes a row to
                    \texttt{history.csv} and calls \texttt{leaderboard.sh}.
              \item Prompts both players to play another game or quit.
          \end{enumerate}
\end{itemize}

\subsection{Game Menu}

The menu renders each game as a plain text row. A single click on any row
launches that game immediately. On hover, a subtle background tint and a thin
left accent bar appear. There are no card boxes, no play button, and no
double-click requirement.

\subsection{In-Game UI}

Every game uses a shared 700$\times$800 window with the following layout:

\begin{center}
\begin{tabular}{|c|c|}
\hline
\textbf{Region} & \textbf{Height} \\
\hline
Top UI bar & 80 px \\
\hline
Board area (centred, with 20 px margin) & 680 px \\
\hline
\end{tabular}
\end{center}

The top bar (implemented in \texttt{base\_game.py} as \texttt{draw\_top\_bar()})
shows:
\begin{itemize}
    \item \textbf{Left}: game name (small, dimmed).
    \item \textbf{Centre}: \textcolor{red}{P1name} \textbf{vs}
          \textcolor{blue}{P2name} in player colours.
    \item \textbf{Right}: whose turn it is (with a coloured dot), or the
          result and keyboard hints once the game ends.
\end{itemize}

% ===================================================================
\section{Class Structure}
% ===================================================================

\subsection{Base Class --- \texttt{BoardGame}}

All games inherit from the \texttt{BoardGame} class defined in
\texttt{game.py}. Each method in the class is overridable
subclass.

\textbf{Shared state attributes:}
\begin{itemize}
    \item \texttt{player\_names} --- dict mapping player integer (1 or 2) to
          username string.
    \item \texttt{current\_player} --- integer (1 or 2) indicating whose turn.
    \item \texttt{board} --- a NumPy array representing the game board.
    \item \texttt{winner} --- \texttt{1}, \texttt{2}, \texttt{0} (draw), or
          \texttt{None} (ongoing).
    \item \texttt{move\_count} --- total moves made.
\end{itemize}

\textbf{Concrete helpers (available to all subclasses):}
\begin{itemize}
    \item \texttt{switch\_turn()} --- flips \texttt{current\_player} between 1
          and 2.
    \item \texttt{is\_game\_over()} --- returns \texttt{True} if
          \texttt{winner} is not \texttt{None}.
    \item \texttt{current\_player\_name()} --- returns the active player's
          username.
    \item \texttt{get\_result\_string()} --- human-readable result for the
          post-game screen.
    \item \texttt{draw\_top\_bar(surf)} --- renders the shared top UI bar.
    \item \texttt{get\_font(size, bold)} --- returns a \texttt{SysFont}.
\end{itemize}

\textbf{Methods subclasses must override:}

\begin{center}
\begin{tabular}{lp{9cm}}
\toprule
\textbf{Method} & \textbf{Purpose} \\
\midrule
\texttt{reset()} & Initialise board to starting state, reset all state. \\
\texttt{check\_winner()} & Return 1, 2, 0, or None using NumPy operations. \\
\texttt{draw\_board(surf)} & Render full frame onto \texttt{surf}; must call
\texttt{draw\_top\_bar()}. \\
\texttt{handle\_click(pos)} & Convert pixel click \texttt{(x,y)} to a game
move. \\
\bottomrule
\end{tabular}
\end{center}

% ===================================================================
\section{Games}
% ===================================================================

\subsection{Tic-Tac-Toe (\texttt{games/tictactoe.py})}

\begin{itemize}
    \item \textbf{Board}: 10$\times$10 NumPy integer array. Values: 0 = empty,
          1 = Player 1 (X), 2 = Player 2 (O).
    \item \textbf{Win condition}: 5 marks in a row --- horizontal, vertical,
          or diagonal.
    \item \textbf{Cell size}: \texttt{best\_cell\_size(10) = 660 // 10 = 66 px}.
\end{itemize}


\subsection{Connect Four (\texttt{games/connect4.py})}

\begin{itemize}
    \item \textbf{Board}: 7$\times$7 NumPy array. Values: 0 = empty, 1 = P1,
          2 = P2.
    \item \textbf{Win condition}: 4 coins in a row.
    \item \textbf{Gravity}: \texttt{np.where(board[:, col] == 0)[0][-1]} finds
          the lowest empty row --- no Python loop.
    \item \textbf{Cell size}: \texttt{best\_cell\_size(7) = 660 // 7 = 94 px}.
    \item \textbf{Coins colours}: one fixed colour per player
          (red for P1, blue for P2); no cycling.
    \item \textbf{Drop preview}: a fully opaque ghost coin is rendered in the
          gap between the top bar and the board, clamped so it is always fully
          visible.
\end{itemize}

\subsubsection*{Win Detection}

Uses the same sliding-consec approach as Tic-Tac-Toe but with
\texttt{WIN\_LEN = 4}. All four directions checked with
\texttt{broadcasting} and \texttt{np.argwhere}.

\subsection{Othello (\texttt{games/othello.py})}

\begin{itemize}
    \item \textbf{Board}: 8$\times$8 NumPy array. Values: 0 = empty,
          1 = Black (P1), 2 = White (P2).
    \item \textbf{Starting position}: two Black and two White discs in the
          centre four cells.
    \item \textbf{Cell size}: \texttt{best\_cell\_size(8) = 660 // 8 = 82 px}.
    \item \textbf{Board appearance}: classic green felt background with corner
          star-point dots.
\end{itemize}


\subsubsection*{Turn Skipping}

After each move, if the next player has no valid moves their turn is skipped
back. If neither player can move, \texttt{check\_winner()} counts discs with
\texttt{np.sum(board == player)} and declares the result.

\subsubsection*{UI Features}

\begin{itemize}
    \item \textbf{Valid move hints}: small dots on all legal cells so players
          always know where they can place.
    \item \textbf{Ghost disc}: semi-transparent disc previews placement on
          hover.
    \item \textbf{Flip ring}: gold ring drawn around discs flipped on the last
          move.
    \item \textbf{Disc count bar}: proportional fill bar below the board
          showing each player's current count.
\end{itemize}

% ===================================================================
\section{Authentication System (\texttt{main.sh})}
% ===================================================================

\begin{itemize}
    \item Uses \texttt{declare} (not \texttt{eval}) for safe variable
          assignment --- prevents code injection from special characters in
          usernames.
    \item Passwords are hashed with:
\begin{lstlisting}[style=bashstyle]
hashed=$(echo -n "$password" | sha256sum | cut -d' ' -f1)
\end{lstlisting}
    \item \texttt{grep -F} with a literal tab separator prevents regex
          false-matches on special-character usernames.
    \item Password confirmation on registration uses its own inner
          \texttt{while true} loop, so a mismatch re-prompts for the password
          rather than looping back to the username prompt.
    \item \texttt{touch "\$USERS\_FILE"} at startup creates the file if it
          does not exist, preventing \texttt{grep} errors on first run.
\end{itemize}

% ===================================================================
\section{Leaderboard (\texttt{leaderboard.sh})}
% ===================================================================

\texttt{leaderboard.sh} is called by \texttt{game.py} after each game, with
the sort metric chosen by the players as an argument
(\texttt{wins}, \texttt{losses}, or \texttt{ratio}).

\begin{itemize}
    \item Reads \texttt{history.csv}, groups results by game name.
    \item For each game, uses \texttt{awk} to accumulate wins and losses per
          player, computes win/loss ratio, and sorts the output using the
          passed metric.
    \item Prints a formatted bordered table to the terminal with columns:
          Rank, Player, Wins, Losses, W/L Ratio.
    \item Handles missing or empty history gracefully with a friendly message.
\end{itemize}



% ===================================================================
\section{Modules Used}
% ===================================================================

\subsection{pygame-ce}
Pygame Community Edition~\cite{pygame} is used for all graphical rendering
across the game hub: the menu screen, all three game boards, the top UI bar,
and the post-game screens. It handles the event loop, keyboard and mouse
input, and font rendering.

\subsection{NumPy}
NumPy~\cite{numpy} is used for all game board representations and win
detection. All boards are stored as NumPy integer arrays.
Win conditions are checked using vectorized operations
(\texttt{np.cumsum}, \texttt{np.where}, \texttt{np.maximum.accumulate},
\texttt{np.diagonal}, \texttt{np.fliplr}) --- no Python loops over cells.

\subsection{Matplotlib}
Matplotlib~\cite{matplotlib} is listed as an allowed module. It is available
for post-game analytics visualisation (e.g.\ bar chart of top players, pie
chart of most-played games from \texttt{history.csv}).

\subsection{pathlib, sys, os}
Standard library modules used for path manipulation, argument parsing, and
cross-platform file handling.

\subsection{subprocess}
Used in \texttt{game.py} to call \texttt{leaderboard.sh} with
\texttt{subprocess.run(["bash", LEADERBOARD\_SH, sort\_by])}.

\subsection{csv, datetime}
Standard library modules used in \texttt{game.py} for writing rows to
\texttt{history.csv} after each game.

% ===================================================================
\section{Directory Structure}
% ===================================================================

\begin{lstlisting}[style=bashstyle]
hub/
|---- main.sh               # Entry point: authentication + launch
|---- game.py               # Game engine: menu, loop, history, leaderboard
|---- leaderboard.sh        # Terminal leaderboard (Bash + awk)
|---- users.tsv             # Username <TAB> SHA-256 hash
|---- history.csv           # Winner, Loser, Date, Game
|---- games/
|    |---- tictactoe.py     # TicTacToe (10x10, 5 in a row)
|    |---- othello.py       # Othello (8x8, flip discs)
|    `---- connect4.py      # ConnectFour (7x7, drop coins)
\end{lstlisting}

% ===================================================================
\section{Basic and Advanced Features}
% ===================================================================

\subsection{Basic Features}

\subsubsection*{Two-Player Authentication}
SHA-256 hashed passwords, registration and login flow, duplicate username
prevention, and persistent storage in \texttt{users.tsv}. Implemented
entirely in Bash.

\subsubsection*{Interactive Game Menu}
Pygame-rendered menu with plain text game rows. Single click launches the
selected game. Shows both player names in their respective colours.

\subsubsection*{Three Playable Games}
All three games are rendered with Pygame GUI windows. No terminal-based
gameplay. All boards are NumPy arrays. All win conditions use NumPy
vectorized operations.

\subsubsection*{Game History Recording}
After each game, \texttt{game.py} appends a row to \texttt{history.csv}
containing Winner, Loser, Date, and Game name. The file is created with a
header row on first use.

\subsubsection*{Terminal Leaderboard}
After each game, the leaderboard is printed to the terminal sorted by the
metric the players chose. Implemented in Bash with \texttt{awk} for per-game
per-player statistics.

\subsection{Advanced Features}

\subsubsection*{Shared Base Class}
\texttt{BoardGame} in \texttt{games/base\_game.py} provides shared constants,
layout helpers, and the \texttt{draw\_top\_bar()} method. All three game
classes inherit from it, keeping game-specific code focused and minimal.

\subsubsection*{Fully Vectorized Win Detection}
Tic-Tac-Toe and Connect Four win detection use zero Python loops for checking.

\subsubsection*{Drop Preview (Connect Four)}
A ghost coin is rendered in the gap between the top bar and board, clamped to
always be fully visible. The column hover highlight is clipped to the board
area.

\subsubsection*{Valid Move Hints and Ghost Disc (Othello)}
All legal placements are shown as dots. A semi-transparent ghost disc
previews placement on hover. Flipped discs receive a gold highlight ring.

\subsubsection*{Post-Game Sort Picker}
The post-game screen is two-phased: first the players pick the leaderboard
sort metric (which triggers CSV write and leaderboard output), then they
choose to play again or quit.

% ===================================================================
\section{Hurdles Faced}
% ===================================================================

\begin{itemize}
    \item \textbf{Vectorizing win detection without sliding window view}:
          The standard approach for detecting runs in a board uses
          \texttt{np.lib.stride\_tricks.sliding\_window\_view}, which was
          not permitted. The solution was the cumsum reset trick, which
          detects consecutive runs across entire rows/columns simultaneously
          using only \texttt{np.cumsum}, \texttt{np.where}, and
          \texttt{np.maximum.accumulate}.

    \item \textbf{Diagonal win detection without loops}:
          Diagonals have varying lengths and cannot be stacked directly.

    \item \textbf{Connect Four drop preview being hidden}:
          The ghost coin was being rendered before \texttt{draw\_top\_bar()},
          which painted over it. Fixed by moving the preview draw call to
          after the top bar, and clamping the coin's \texttt{y} coordinate
          to the gap between the bar and board.

    \item \textbf{Avoiding circular imports}:
          Game files need shared constants and the base class, but
          \texttt{game.py} dynamically imports game files. Placing
          \texttt{BoardGame} in \texttt{game.py} would create a circular
          import. The solution is \texttt{games/base\_game.py}: a pure data
          and class file that imports nothing from the rest of the hub.

    \item \textbf{Othello turn-skip logic}:
          After placing, if the next player has no valid moves, the turn must
          skip back. If neither can move, the game ends. This required careful
          ordering: switch turn $\to$ check valid moves $\to$ possibly switch
          back $\to$ check valid moves again $\to$ end game.
\end{itemize}

% ===================================================================
\section{What I Would Do Differently With 10 More Hours}
% ===================================================================

\begin{itemize}
    \item \textbf{Fullscreen support}: Pygame's
          \texttt{pygame.display.toggle\_fullscreen()} is trivial to add, but
          it only scales the window without scaling the board. A proper
          implementation would require making all layout constants dynamic
          (computed from live screen size at runtime) rather than module-level
          constants. With more time, this refactor would be worthwhile.

    \item \textbf{Matplotlib analytics}: The project specification requires
          a bar chart of top 5 players and a pie chart of most-played games.
          These are straightforward with \texttt{history.csv} and Matplotlib
          but were not implemented due to time constraints.

    \item \textbf{Animated coin drop in Connect Four}: Rather than the coin
          instantly appearing in its final row, a smooth fall animation along
          the column would improve the game feel.


    \item \textbf{Sound effects}: Short Pygame mixer sounds for placing marks,
          flipping discs, and winning would add polish.

    \item \textbf{Coin flip animation in Othello}: Smoothly transitioning disc
          colour when flipped (via a squash/stretch animation) would make the
          game more engaging.
\end{itemize}

% ===================================================================
\section{References}

\begin{enumerate}
    \item Python Standard Library Docs: \\
    \texttt{https://docs.python.org/3/library}

    \item Matplotlib Documentation: \\
    \texttt{https://matplotlib.org/}


    \item Matplotlib Tick Locators: \\
    \texttt{https://matplotlib.org/stable/gallery/ticks/tick-locators.html}

    \item Matplotlib Tick Formatters: \\
    \texttt{https://matplotlib.org/stable/gallery/ticks/tick-formatters.html}
\end{enumerate}

\end{document}
